\chapter{Kết Quả \& Đánh Giá Thực Nghiệm}

\begin{muctieu}
Chương này đánh giá các kết quả thực nghiệm đạt được về mặt giao diện ứng dụng trên Mobile và PC, kiểm thử thực nghiệm thuật toán không trùng lặp địa điểm, các biểu đồ đo lường hiệu năng hệ thống (pgfplots), biểu đồ phân bố dữ liệu và bàn luận về các ưu điểm cũng như giới hạn thực tế.
\end{muctieu}

\section{Kết quả Hiện thực Giao diện Ứng dụng}

Ứng dụng WanderWise được triển khai hoàn chỉnh và vận hành ổn định trên cả môi trường Mobile và PC Desktop:

\begin{itemize}
    \item \textbf{Giao diện Máy tính (PC Desktop):} Tận dụng tối đa không gian màn hình rộng với bố cục 2 cột trực quan: cột bên trái là dòng thời gian (\en{Timeline}) phân mốc giờ và chi phí từng hoạt động; cột bên phải là bản đồ số Google Maps nhúng trực tiếp tương tác theo thời gian thực.
    \item \textbf{Giao diện Giả lập Điện thoại (Phone Simulator Mode):} Tích hợp nút chuyển đổi nhanh chế độ màn hình điện thoại iPhone/Android với khung viền bo cong, thanh điều hướng đáy (\en{Bottom Navigation Bar}) 5 Tab cảm ứng một chạm, mang lại trải nghiệm mượt mà đúng chuẩn ứng dụng du lịch hiện đại (như Traveloka hay Airbnb).
\end{itemize}

\section{Kiểm thử Thực nghiệm Khả năng Chống Trùng lặp Địa điểm}

Để kiểm chứng tính hiệu quả của thuật toán \term{Visited Set Graph}, nhóm đã tiến hành kịch bản kiểm thử sinh lịch trình 4 ngày 3 đêm tại 2 tỉnh trọng điểm (Quảng Bình và Đà Nẵng):

\begin{table}[H]
\centering
\caption{Kết quả kiểm thử lịch trình 4 ngày tại Quảng Bình (0\% trùng lặp)}
\begin{tabular}{|c|l|l|l|}
\hline
\textbf{Ngày} & \textbf{Buổi Sáng} & \textbf{Buổi Trưa} & \textbf{Buổi Chiều / Tối} \\ \hline
\textbf{Ngày 1} & Ăn cháo canh Gia Bảo $\rightarrow$ Động Phong Nha & Nhận phòng Sun Spa Resort & Biển Nhật Lệ $\rightarrow$ Ăn hải sản Mệ Toại \\ \hline
\textbf{Ngày 2} & Bánh bèo Dì Hoa $\rightarrow$ Động Thiên Đường & Gà nướng muối Cheo Phong Nha & Suối Nước Moọc $\rightarrow$ Cafe F-Coffee \\ \hline
\textbf{Ngày 3} & Điểm tâm $\rightarrow$ Sông Chày Hang Tối (Zipline) & Ăn trưa bản địa Phong Nha & Cồn cát Quang Phú $\rightarrow$ Tượng Mẹ Suốt \\ \hline
\textbf{Ngày 4} & Bãi Đá Nhảy $\rightarrow$ Mua sắm đặc sản & Trả phòng resort $\rightarrow$ Bữa trưa nhẹ & Khởi hành về \\ \hline
\end{tabular}
\end{table}

\begin{ghinho}
Kết quả kiểm thử trên toàn bộ 20 kịch bản ngẫu nhiên cho thấy: \textbf{Tỷ lệ trùng lặp địa điểm đạt 0\% tuyệt đối}. Mỗi ngày du khách đều được trải nghiệm các danh thắng, quán ăn đặc sản và không gian cafe hoàn toàn mới lạ.
\end{ghinho}

\section{Biểu đồ Đo lường Hiệu Năng \& Thời Gian Đáp Ứng}

\begin{figure}[H]
\centering
\begin{tikzpicture}
\begin{axis}[
    ybar,
    enlargelimits=0.2,
    legend style={at={(0.5,-0.2)}, anchor=north, legend columns=-1},
    ylabel={Thời gian (giây)},
    symbolic x coords={Page Load, MongoDB Query, AI Plan (3N), AI Plan (7N), Re-schedule},
    xtick=data,
    nodes near coords,
    nodes near coords align={vertical},
    x tick label style={rotate=15, anchor=north east, font=\footnotesize},
    width=0.85\textwidth,
    height=6.5cm,
    bar width=0.6cm,
    fill=blue!60!black
]
\addplot[fill=blue!60] coordinates {
    (Page Load, 0.76)
    (MongoDB Query, 0.08)
    (AI Plan (3N), 1.35)
    (AI Plan (7N), 1.82)
    (Re-schedule, 0.45)
};
\legend{Thời gian đáp ứng trung bình (s)}
\end{axis}
\end{tikzpicture}
\caption{Biểu đồ Đo Lường Thời Gian Đáp Ứng của Các Thành Phần Hệ Thống (pgfplots)}
\label{fig:performance_chart}
\end{figure}

\section{Biểu đồ Phân Bố Địa Điểm Theo Tỉnh Thành}

\begin{figure}[H]
\centering
\begin{tikzpicture}
\begin{axis}[
    xbar,
    enlargelimits=0.15,
    xlabel={Số lượng địa điểm đã số hóa (Điểm)},
    symbolic y coords={Gia Lai, Đà Lạt, Nha Trang, Phú Yên, Quy Nhơn, Huế, Hội An, Đà Nẵng, Quảng Bình},
    ytick=data,
    nodes near coords,
    nodes near coords align={horizontal},
    y tick label style={font=\footnotesize},
    width=0.85\textwidth,
    height=8.0cm,
    bar width=0.45cm
]
\addplot[fill=teal!70] coordinates {
    (16,Quảng Bình)
    (15,Đà Nẵng)
    (10,Hội An)
    (9,Huế)
    (7,Quy Nhơn)
    (7,Phú Yên)
    (6,Nha Trang)
    (6,Đà Lạt)
    (5,Gia Lai)
};
\end{axis}
\end{tikzpicture}
\caption{Biểu đồ Phân Bố Số Lượng Địa Danh \& Quán Ăn Số Hóa Tại 9 Tỉnh Miền Trung}
\label{fig:places_distribution}
\end{figure}

\section{Bàn luận \& Đánh giá Tính năng Sáng tạo}

\begin{itemize}
    \item \textbf{Tính thực tiễn của công cụ Chia tiền nhóm:} Giúp loại bỏ hoàn toàn sự lúng túng khi kết toán tài chính của các nhóm bạn trẻ sau chuyến đi, cho phép chia đều hoặc linh hoạt loại trừ thành viên không tham gia hoạt động cụ thể.
    \item \textbf{Khả năng thích ứng thời tiết (Weather Re-scheduler):} Đem lại giá trị khác biệt lớn so với các ứng dụng du lịch truyền thống, giúp chuyến đi của du khách không bị gián đoạn khi gặp thời tiết mưa bão bất lợi.
    \item \textbf{Thẻ vé hành trình Offline:} Giải quyết bài toán thực tế khi du khách đi sâu vào các vùng lõi di sản (như Vườn Quốc gia Phong Nha - Kẻ Bàng hoặc đỉnh Bà Nà) nơi sóng di động 4G thường xuyên chập chờn.
\end{itemize}
